% Source-aligned English for printed p.149.
\SourcePage{149}

\LayerTarget{tohoku:section:2.5}
\textbf{2.5. Resolving functors.} 
\LayerTarget{tohoku:def:resolving-functor}
Let $\mathbf{C}$ and $\mathbf{C}'$ be two abelian categories, and suppose every object of $\mathbf{C}$
is isomorphic to a subobject of an injective object. Let $F$ be a covariant
left exact functor from $\mathbf{C}$ to $\mathbf{C}'$. A \textit{resolving functor} for $F$ is
a covariant functor $\mathbf F$ on $\mathbf{C}$ with values in the category of
nonnegatively graded complexes in $\mathbf{C}'$, with
$\mathbf F(A)=(\mathbf F^n(A))$, equipped with an augmentation
$F\to\mathbf F$. Thus there is a homomorphism from $F$ to the
functor $Z^0\mathbf F^0$ of cocycles of $\mathbf F^0$. It must satisfy:
(i) $\mathbf F$ is exact; (ii) $F\to Z^0\mathbf F^0$ is an
isomorphism; and (iii) if $A$ is injective, then $\mathbf F(A)$ is acyclic in
degrees $>0$.

Let $\mathbf F$ be a resolving functor for $F$. Consider the functors
$H^n\mathbf F(A)$ on $\mathbf{C}$ with values in $\mathbf{C}'$. By (i) they form a
cohomological functor; by (ii) this functor reduces to $F$ in degree 0; and by
(iii) the $H^n\mathbf F(A)$ are effaceable for $n>0$. Consequently:

\LayerTarget{tohoku:prop:2.5.1}
\textsc{Proposition 2.5.1.} \textit{If $\mathbf F$ is a resolving
functor for $F$, then for every $A\in \mathbf{C}$ there are unique isomorphisms
$H^n\mathbf F(A)=\mathbf{R}^nF(A)$. They define an isomorphism of cohomological
functors and reduce in degree $0$ to the augmentation isomorphism.}

We shall give another proof of the preceding proposition, which provides a
convenient computation of these isomorphisms.

\LayerTarget{tohoku:prop:2.5.2}
\textsc{Proposition 2.5.2.} \textit{Let $\mathbf F$ be a resolving
functor for $F$. Let $A\in \mathbf{C}$, and let $C=(C^p(A))$ be a right resolution of
$A$ by $F$-acyclic objects. Consider the bicomplex
$\mathbf F C(A)=(\mathbf F^qC^p(A))_{p,q}$ and the natural homomorphisms}
\LayerTarget{tohoku:eq:2.5.1}
\[
\tag{2.5.1}
F(C(A))\to\mathbf F(C(A))\leftarrow\mathbf F(A)
\]
\textit{defined respectively by the augmentation
$F\to\mathbf F$ and the augmentation
$A\to C(A)$. Then the corresponding homomorphisms}
\LayerTarget{tohoku:eq:2.5.2}
\[
\tag{2.5.2}
H^nF(C(A))=R^nF(A)\to H^n\mathbf F(C(A))
\leftarrow H^n\mathbf F(A)
\]
\textit{are isomorphisms, and the corresponding isomorphism
$H^n\mathbf F(A)=R^nF(A)$ is the one in Proposition \textup{2.5.1.}}

\LayerTarget{tohoku:prop:2.5.2:proof:p1}
Regard $F(C(A))$, respectively $\mathbf F(A)$, as a bicomplex whose second,
respectively first, degree is zero. The first spectral sequence of
$\mathbf F(C(A))$ has initial term $H(F(C(A)))$, because
$H_{II}^q(\mathbf F(C(A)))$ vanishes for $q>0$ and identifies with
$F(C(A))$ for $q=0$, by Proposition 2.5.1 and because the $C^p(A)$ are
$F$-acyclic. The first homomorphism in (2.5.1) therefore induces an
isomorphism on the initial terms of spectral sequence $I$, hence an
isomorphism on cohomology. Similarly,
$H_I(\mathbf F(C(A)))^{p,q}$ vanishes for $p>0$ and equals
$\mathbf F^q(A)$ for $p=0$, since $\mathbf F$ is exact and $C(A)$ is a
resolution of $A$. Thus the second homomorphism in (2.5.1) induces an
isomorphism on the initial terms of spectral sequences $II$, and hence again
on cohomology. To show that the resulting isomorphism from
$H^n\mathbf F(A)$ to $R^nF(A)$ is the one in Proposition 2.5.1, take an
\textit{injective} resolution $C'(A)$ of $A$, a homomorphism from $C(A)$ to $C'(A)$,
and consider the associated homomorphism of diagram (2.5.2)

\TohokuContinuation{150}
